\begin{frame}{동적계획법의 핵심 아이디어}
  \begin{block}{``큰 문제 = 작은 부분 문제 + 재귀''}
    이 패턴은 \textbf{컴퓨터공학 전반}에서 반복해서 등장한다.
    강화학습은 그 패턴을 ``미래의 \textbf{불확실한} 보상을
    어떻게 \textbf{지금} 계산할 것인가''라는 문제에 적용한 것.
  \end{block}
  \vskip0.5em
  \begin{itemize}
    \item 재귀 구조(벨만방정식)로 \textbf{무한한 미래}를
      \textbf{한 스텝}으로 압축.
    \item 그 반복이 \textbf{항상 정답에 수렴}한다는 것 자체가,
      \kb{축소 사상(contraction mapping)}이라는 \textbf{하나의
      수학적 성질}로 \textbf{보장}된다.
  \end{itemize}
  \vskip0.5em
  {\footnotesize 표준 참고: Sutton \& Barto(2018) 4장
  ``Dynamic Programming'' / David Silver UCL 강화학습 강의 /
  Stanford CS234. 4.1$\sim$4.3절의 모든 알고리즘은 ``재귀로
  큰 문제를 쪼개는'' 같은 패턴의 세 가지 변주(기대/최적 $\times$
  V/Q)일 뿐이다.}
\end{frame}
